\documentclass[letterpaper, 12pt]{article}
\usepackage[top=2 cm, bottom=2 cm,left=1.5 cm, right=1.5 cm]{geometry}

\usepackage[spanish]{babel}
\usepackage[utf8]{inputenc}
\usepackage{csquotes}

\setlength{\parindent}{1cm}
\pagestyle{empty}

\usepackage{amsmath}
\usepackage[outputdir=aux]{minted}
\usepackage{amssymb}
\usepackage{multicol}
\usepackage{enumerate}
\usepackage{graphicx} % Incluir imágenes
\usepackage[document]{ragged2e}
\usepackage{mathrsfs}
\usepackage[T1]{fontenc}
\usepackage[table,xcdraw]{xcolor}
\usepackage{float} %required for the placement specifier H
\makeatletter
\makeatother
% -----------------------------------------------

% Preámbulo

\title{Tarea}
\author{Christopher Gómez}
\date{Marzo 2022}

\begin{document}
\parbox[t]{.5\linewidth}{
    \centering
    \includegraphics[scale=0.4]{logo.png}
    \begin{center}
        UNIVERSIDAD SIMÓN BOLÍVAR \\
        CI-5651 - Diseño de Algoritmos I \\
        Prof. Ricardo Monascal \\
    \end{center}
}
\hfill \framebox[5.5cm][c]{
        \parbox[t]{.8\linewidth}{
        \centering
       Resuelto por:\\ Christopher Gómez
        }}

\phantom{This text will be invisible} \\
\centerline {\textbf{Tarea}}
\justify

\begin{enumerate}

% -------- PREGUNTA 1 --------

\item Considere el algoritmo \texttt{stupidSort} para ordenar un arreglo de menor a
mayor: \\

\begin{minted}{python}
def stupidSort(a):
    for x in permutaciones(a):
        if ordenado(x):
            return x
\end{minted}

Puede suponer que \texttt{permutaciones} es un iterador que devuelve, una a una, todas las permutaciones de un arreglo en tiempo constante. Además, \texttt{ordenado} es un predicado que verifica si un arreglo está ordenado de menor a mayor en tiempo $\Theta(n)$. \\

Diga cotas inferior y superior para \texttt{stupidSort} (correspondiendo al peor y mejor caso, respectivamente). Esto es, defina dos funciones $f$ y $g$ tal que se cumpla que: \\

\begin{multicols}{2}
    \begin{itemize}
        \item \texttt{stupidSort} $\in O(f)$
        \item \texttt{stupidSort} $\in \Omega(g)$
    \end{itemize}
\end{multicols}

Para el mejor caso, supongamos que el primer elemento que retorna el iterador \texttt{permutaciones} es el arreglo ordenado. En este caso, el algoritmo \texttt{stupidSort} retorna el arreglo ordenado en tiempo $\Theta(n)$, que es el tiempo que tarda \texttt{ordenado} en verificar que el arreglo está ordenado. Por lo tanto, $g(n) = n$. \\

Por otro lado, el peor caso se da cuando el arreglo ordenado resulta ser el último retornado por {permutaciones}, lo cual solo ocurre cuando el iterador ha agotado todas las permutaciones posibles sobre un arreglo de $n$ elementos, que son $n!$. Para cada una de estas permutaciones \texttt{ordenado} verifica si está ordenado en $\Theta(n)$, por lo que el tiempo total es $O(n! \times n)$. Así, $f(n) = n! \times n$. \\

\item  En la sede del banco caben hasta C personas. La sede empieza vacía y los
clientes entran uno a uno. Cuando la capacidad de la sede es superada, la misma cierra y la totalidad de los clientes es transferida a una sede más grande, con capacidad $2 \times C$.

Nótese que si contamos cuántas personas deben trasladarse al entrar un cliente a la sede, ocurre uno de dos casos:

\begin{itemize}
    \item Si aún no se ha superado la capacidad de la sede, la cantidad de personas que debe trasladarse es $O(1)$ — únicamente el cliente que entra.
    \item Si se supera la capacidad de la sede, la cantidad de personas que debe trasladarse es $O(C)$ — todos los clientes en la sede del banco más el cliente nuevo.
\end{itemize}

Demuestre que el orden amortizado, en el peor caso, para el ingreso de un cliente es $O(1)$. \\

\item El problema $\oplus-$SAT modificación sobre el problema de satisfacibilidad booleana, donde la fórmula proposicional está en forma normal conjuntiva, pero cada cláusula está separada por disyunciones exclusivas (en vez de disyunciones tradicionales).

Recordemos que $P \oplus Q$ es cierto si y solo si uno de entre $P$ y $Q$ es cierto (el otro siendo falso).

Diga si $\oplus-$SAT $\in P$ o si $\oplus-$SAT es $NP-$completo. Demuestre su afirmación.



\end{enumerate} \vspace{4mm}

\end{document}
